\section{Reviewer 1}

\subsection*{R1-1 --- Prominence of the isotropic-penalty comparison, and whether the Fisher mechanism earns anything}

\begin{rquote}
The new experiments show that Fisher weighting does not improve the best attainable accuracy
relative to a separately tuned equal-weight penalty. More importantly, a simple isotropic
displacement penalty achieves essentially the same accuracy as FOLoRA. The authors should
give this comparison greater prominence and establish whether the Fisher-based mechanism
provides a meaningful benefit in forgetting, computational efficiency, or another
independently evaluated property.
\end{rquote}

\resp We accept both halves of this. The comparison is now in the main results table rather
than in the limitations list, and we looked for a benefit of the Fisher mechanism on each of
the three axes you name. We found none that we can defend, and the paper now says so; the
claims that depended on the Fisher weighting being the source of the margin are gone from the
title, abstract and highlights (R8-3).

\paragraph{Prominence.} The isotropic penalty is now a row of the main comparison,
Table~1, immediately above the FOLoRA row: a shared adapter trained with
$\lambda\sum_i\lVert\delta\theta_i\rVert^2$ scaled to the same RMS as the Fisher term and
evaluated on the same three class orders. Its entries are $78.27\pm0.38$ on CIFAR-100 at the
matched strength $\lambda{=}3$ ($77.77\pm0.56$ at $\lambda{=}10$) against FOLoRA's
$78.37\pm0.84$ over the same three seeds; on ImageNet-R the same three seeds give
$66.89\pm0.29$ against FOLoRA's $66.89\pm0.41$ at $\lambda{=}3$ (paired $\Delta{=}+0.00$,
$p{=}0.98$).

\paragraph{Forgetting.} On the gap metric (CIL$-$TIL: how far accuracy drops when the task
identity is removed) the contrast is inside the noise in both directions: on CIFAR-100 the
isotropic penalty is nominally \emph{better} ($6.83\pm0.35$ against $7.05$ at
$\lambda{=}3$, paired $\Delta{=}-0.22$, $p{=}0.49$; $6.71\pm0.13$ against $6.88$ at
$\lambda{=}10$, $\Delta{=}-0.17$, $p{=}0.28$), and on ImageNet-R nominally \emph{worse}
($7.30\pm0.14$ against $7.11\pm0.42$, $\Delta{=}+0.18$, $p{=}0.38$). Neither contrast is
significant, and the sign is not consistent across the two benchmarks, so we do not report the
gap as evidence for the Fisher mechanism in either direction.

\paragraph{Computational efficiency.} The isotropic arm is not slower in our implementation
--- it shares the same per-sample Fisher estimator, so training cost is identical --- but it is
strictly cheaper to \emph{store}: one scalar per layer, $\sum_j\sigma_j^2/d$, in place of the
rank-$k$ basis that FOLoRA retains ($k\times d$ per layer, $182.8$\,MiB in total at $k{=}64$).
A standalone L2-SP implementation would not need the Fisher estimate at all and
would be cheaper still; our run therefore \emph{understates} the efficiency advantage of the
simpler penalty. On this axis the comparison favours the isotropic penalty, and
Sec.~5.2 now says that a practitioner who wants only the accuracy should
prefer it.

\paragraph{Another independently evaluated property.} We could not establish one. The rank-$k$
Fisher basis constrains only the directions past tasks were measured to be sensitive to and
leaves the remaining $d-k$ dimensions free, while the isotropic penalty constrains all of
them; that is a qualitative difference in what is protected, and at our budget, task count
and backbone it does not translate into a difference in accuracy, forgetting, or the
decision-rule audit. The experiments that could have separated them (equal-weight,
random-direction and isotropic penalties at matched scale and matched $\lambda$) all land
within one point of FOLoRA; at their own best strengths the weighted and equal-weight arms
are indistinguishable ($+0.43$, $p=0.44$, five paired seeds; Sec.~5.5), and the TOST
analysis in Sec.~5.6 --- run for the one comparison whose conclusion is parity, FOLoRA
against tuned EWC-LoRA --- reports equivalence at $\delta{=}1.0$.

\paragraph{What we changed.} (i) The isotropic comparison is a main-table row and a paragraph
of Sec.~5.2, not only limitation (vii). (ii) The title now names what the
evidence supports --- one shared adapter rather than one module per task --- instead of
leaning on the Fisher adjective; the abstract states that the margin over the
parameter-efficient family is attributable to regularizing a shared adapter and that the
Fisher weighting is not the source, and the highlights say the same (R8-3). (iii) The
contribution list in Sec.~1 distinguishes what is a \emph{method} (a shared-adapter
regularizer with a matrix-valued weight) from what is a \emph{finding} (the choice of weight
matters much less than the fact that the adapter is shared), and the abstract reports the
latter as a finding rather than implying the former is the source of the gain.

\change Table~1 (new isotropic row), Sec.~5.2 (new
paragraph), Sec.~5.6 limitation (vii) rewritten to point at the table
instead of carrying the only copy of the comparison, Sec.~5.1 (the baseline list now
enumerates the isotropic row and the seed-budget list gives its $3$/$3$), title, abstract,
highlights, Sec.~1 contribution list, conclusion.

\subsection*{R1-2 --- The main table now runs O-LoRA at its training-time objective}

\begin{rquote}
The principal results table still uses an initialization-only O-LoRA variant, even though the
authors have evaluated the original training-time orthogonality objective and obtained
substantially stronger results. The faithful configurations should appear in the main
comparison, with the initialization-only version labeled as an ablation.
\end{rquote}

\resp Accepted, and it was the right call: reporting the initialization-only variant against
baselines at their full objective was not a like-for-like comparison, and the variant we had
in the table was also the only row below the frozen-feature floor on CIFAR-100.

\paragraph{What changed.} Table~1 now reports O-LoRA at
$\lambda_1\sum_{i<t}\lVert A_t^{\top}A_i\rVert_F^2$ with $\lambda_1{=}1.0$ --- the stronger
of the two strengths we ran, and the objective the original paper specifies --- on
CIFAR-100 ($n{=}5$) and ImageNet-R ($n{=}3$). The initialization-only variant is now an
ablation row of Table~A.10, labelled as such, together with the weaker
training-time strength $\lambda_1{=}0.1$ so that the effect of the constraint is visible.
Table~A.10 carries each configuration's paired difference and $p$-value; the accuracies,
which Sec.~5.2 quotes (the training-time row's also being Table~1's), are:

\begin{center}\small
\begin{tabular}{@{}lcc@{\hspace{2em}}cc@{}}
\toprule
& \multicolumn{2}{c}{CIFAR-100} & \multicolumn{2}{c}{ImageNet-R} \\
\cmidrule(lr){2-3}\cmidrule(lr){4-5}
O-LoRA configuration & ACC & GAP & ACC & GAP \\
\midrule
initialization only (now ablation) & 64.94 & --- & 53.72 & --- \\
training-time, $\lambda_1{=}0.1$ & $71.21\pm1.10$ & --- & 56.30 & --- \\
training-time, $\lambda_1{=}1.0$ (now main) & $72.72\pm0.90$ & $9.09\pm0.37$ & $61.15\pm0.92$ & $8.86\pm0.33$ \\
\bottomrule
\end{tabular}
\end{center}

\paragraph{The consequence for our own claims, which we have not softened.} The margin over
O-LoRA falls from $+12.92$ to $+5.60$ on CIFAR-100 and from $+13.20$ to $+5.74$ on
ImageNet-R (paired, common seeds). On CIFAR-100 the difference remains significant after
Holm correction (accuracy $p_{\mathrm{Holm}}{=}0.0015$; gap $p_{\mathrm{Holm}}{=}0.0029$). On
ImageNet-R the accuracy difference no longer clears the corrected threshold (uncorrected
$p{=}0.017$, $p_{\mathrm{Holm}}{=}0.052$); the gap difference does ($0.011$, the value
Sec.~5.1 prints). The abstract,
Sec.~5.2 and the conclusion now state the ImageNet-R accuracy comparison as
just outside the threshold rather than as an improvement. The initialization-only row's
comfortable margin was an artefact of the weaker baseline, and quoting it would have
overstated the result; the paper no longer does.

\paragraph{The comparison is now against baselines at their own budgets.} Because the family
of confirmatory tests changed, every Holm-corrected $p$-value was recomputed rather than
reported from the previous family. The recalculation is mechanical
(\texttt{scripts/significance.py}, family $m{=}9$ on CIFAR-100 and $m{=}8$ on ImageNet-R, not
pooled) and the paper's numbers regenerate byte-for-byte from the per-run JSON on disk.

\paragraph{Collateral changes.} Changing this row touches more than the row:
\begin{itemize}
\item The caption of Table~1 now states the O-LoRA seed budget ($5$/$3$) and the
  objective, so the reader can see that the two methods with the largest seed budgets are
  EWC-LoRA and FOLoRA.
\item Table~A.10 gained the two training-time O-LoRA rows; the
  initialization-only rows remain, relabelled as the ablation, so that both configurations
  are visible and neither is deleted.
\item Table~2 carries the same checkpoints, so its O-LoRA row moves with the main row: head
  accuracy $7.47\to8.71$ and NCM $64.94\to72.72$ on CIFAR-100, and $6.64\to6.58$ and
  $53.72\to61.15$ on ImageNet-R.
\item The Holm family for each benchmark was rebuilt from the new row
  (\texttt{scripts/significance.py}), and every corrected $p$-value that depends on it was
  updated together: in Sec.~5.2, the abstract and the conclusion; in the caption of
  Table~3, whose two ImageNet-R values become $0.0013$ and $0.0032$; and in Table~A.9,
  where the two O-LoRA rows change with the row itself (their budgets and differences
  become $5$/$+5.60$ and $3$/$+5.74$, and neither corrected value is $0.0000$ any longer)
  and six further entries move with the family.
\item The floor discussion in Sec.~5.2 was rewritten: the
  initialization-only O-LoRA was the row that fell below the CIFAR-100 floor ($64.94$ against
  $70.31$); at the training-time objective it clears the floor ($72.72$), so the text now
  names only Seq-LoRA as below the floor and keeps the initialization-only variant's shortfall
  as a further reason it belongs in the ablation.
\item The seed-variance paragraph in Sec.~5.6 that attributed a
  $\pm3.03$ spread to O-LoRA was updated to name the initialization-only variant, since that
  spread does not belong to the row now in the main table.
\item Figures~D.1 and D.2 (candidate-set curves) were recomputed from the new O-LoRA
  checkpoint. Its curve now ends at $72.72$ (CIFAR-100, five seeds) and $61.15$
  (ImageNet-R, three seeds) at the full candidate set, where the initialization-only
  checkpoint ended at $64.94$ and $53.72$. The change is visible in the plotted range: on
  CIFAR-100 the lowest curve is no longer O-LoRA but Seq-LoRA, so the vertical axis of that
  figure is now bounded at $65$ instead of $60$, and the caption and the sentence in
  Sec.~5.2 that quote the plotted spread were updated to the recomputed
  values ($95.1$--$98.2$ at the smallest candidate set and $69.3$--$78.0$ at the full one,
  against the previous $93.6$--$98.2$ and $64.9$--$78.0$); on ImageNet-R the same
  substitution moves the lowest plotted value from $53.72$ to $56.58$, so the caption's
  margin above the axis floor changes from $3.72$ to $6.58$ points. The captions also now
  state the seed budget each curve is drawn on ($5$ for O-LoRA on CIFAR-100, $3$ on
  ImageNet-R) and the isotropic row among those tabulated but not drawn.
\item The RanPAC column of Table~4 was recomputed for the new O-LoRA checkpoint. The
  recomputation covers five class orders on CIFAR-100 and three on ImageNet-R, so the
  cells of the other rows move by a few tenths as well (FOLoRA $84.98\to85.17$, EWC-LoRA
  $85.04\to85.25$), and every RanPAC number quoted in Sec.~5.3 was updated with them
  rather than only the O-LoRA row. For O-LoRA the cell was empty in the submitted version
  and is now $83.20\pm0.77$ on CIFAR-100 and $68.15\pm0.90$ on ImageNet-R. FOLoRA remains
  ahead of it under this rule by $+1.97$ ($p=0.0025$, $n=5$) and $+3.59$ ($p=0.046$,
  $n=3$), and the parity with the tuned EWC-LoRA is unchanged ($85.17$ against $85.25$,
  $p=0.70$). The NCM column of that table continues to come from the same checkpoint as
  the RanPAC column, so the two columns remain a single model.
\item The abstract, highlights and conclusion no longer list O-LoRA among the ImageNet-R
  accuracy wins; the sentence now separates the CIFAR-100 win from the ImageNet-R
  near-threshold result.
\end{itemize}

\change Table~1; Table~2; Table~3 (caption); Table~A.9; Table~A.10;
Sec.~5.1 (seed budgets and the Holm-family paragraph); Sec.~5.2; Sec.~5.3; Sec.~5.6;
Figs.~D.1 and D.2; abstract; highlights; conclusion.
